\section{Introduction}
\label{sec:intro}

Point-of-Interest (POI) prediction in Location-Based Social Networks couples two tasks over a shared substrate. Given a length-9 window of check-ins, we predict the next visit's semantic category and spatial region. The tasks place disjoint demands: \textbf{category prediction} is a 7-way decision shaped by per-visit context (intent, time-of-day, co-location); \textbf{region prediction} is a fine-grained ranking over 1.1K to 8.5K classes hinging on spatial-temporal regularities and short-horizon transitions. The MTL appeal is standard: a shared representation should let one head's signal transfer and lift both~\cite{caruana1997multitask,vandenhende2022mtl}.

Two recent works frame the bottleneck. The first reported hard-parameter-sharing MTL on POI-stable graph embeddings (DGI~\cite{velickovic2019dgi}) yielding only marginal gains within standard deviations; the diagnosis was representation mismatch~\cite{silva2025mtlnet}. A follow-up decomposed the monolithic embedding into spatial, temporal, and categorical sub-encoders, recovering category-side gains across three states~\cite{paiva2026courb}. Together these suggest embedding choice, not the MTL recipe, is the load-bearing factor.

Both prior responses treat the embedding as a \emph{what} axis. We test an orthogonal \emph{granularity} axis: emission rate. POI-stable embeddings like HGI~\cite{huang2023hgi} emit one vector per place; a check-in-level variant emits one per visit, capturing intent shifts across two visits to the same caf\'e at different hours. Our hypothesis is that \textbf{per-visit context} is the load-bearing substrate property for the easier head, and that once the substrate is fixed, joint training behaves like the textbook MTL tradeoff. The carry is task-asymmetric: per-visit \emph{carries} next-category; per-place embeddings tie or marginally exceed it on next-region.

We adopt Check2HGI (a check-in-level extension of the hierarchical graph-infomax line emitting 64-dim per-visit vectors) and run a controlled ablation on five Gowalla~\cite{cho2011gowalla} U.S.-state splits (AL, AZ, FL, CA, TX) under user-disjoint folds. We compare Check2HGI against HGI under matched-head STL baselines, and run a cross-attention MTL backbone (per-task input streams, GRU category head, STAN-Flow region head~\cite{luo2021stan}). The contribution is empirical and methodological: matched-head STL ceilings let us separate substrate effects from head-capacity effects, a per-visit counterfactual decomposes the substrate gain itself, and an encoder-frozen isolation protocol surfaces a cross-attention pitfall in loss-side ablations.

The main contributions of this paper are as follows:
\begin{itemize}
    \item \textbf{C1. A controlled five-state matched-head ablation that separates substrate from head capacity.} Under the matched-head ceiling, check-in-level Check2HGI lifts single-task category macro-F1 over POI-stable HGI by $+14.5$ to $+29$ percentage points (pp) at every state (paired Wilcoxon $p = 0.0312$ each; 5/5 folds positive; head-invariant at AL/AZ across linear, GRU, single-layer, and LSTM probes, and matched-head replicated at FL/CA/TX), while on next-region per-place embeddings tie or marginally exceed Check2HGI (TOST non-inferiority at $\delta = 2$ pp passes at CA and TX). A pooled-vs-canonical counterfactual at Alabama attributes about $72\,\%$ of the category substrate gap to per-visit context, single-state mechanism evidence consistent with the asymmetry.
    \item \textbf{C2. A direct measurement of the multi-task tradeoff under the load-bearing substrate.} With Check2HGI fixed, joint cross-attention MTL adds a small category lift at four of five states (AZ $+1.20$, FL $+1.40$, CA $+1.68$, TX $+1.89$ pp at $p \leq 2\mathrm{e}{-6}$ across $n = 20$ multi-seed fold-pairs) and is small-significantly negative at Alabama ($\Delta = -0.78$ pp, $p = 0.036$, $n = 20$; magnitude under $2\,\%$ relative). On next-region it pays a sign-consistent $7$ to $17$ pp cost on Acc@10 at every state ($p \leq 1.9\mathrm{e}{-6}$, $n = 20$); drop-in alternatives (FAMO, Aligned-MTL, a hierarchical-additive-softmax region head) do not close the gap, indicating that the cost is structural to the cross-attention shared-backbone interaction rather than a recipe artefact.
    \item \textbf{An encoder-frozen isolation protocol for cross-attention MTL} (developed in Section~\ref{sec:mechanism}). A loss-side \texttt{task\_weight = 0} ablation does not silence the targeted encoder, since it co-adapts through the attention key and value projections of the surviving stream; freezing the encoder at random initialisation cleanly decomposes the architectural cost, backed by regression tests in our anonymous code release.
\end{itemize}

Sections~\ref{sec:related}, \ref{sec:method}, and~\ref{sec:setup} cover related work, method, and setup. Section~\ref{sec:results} reports substrate ablation, MTL-vs-STL, $\Delta_m$, and external baselines. Section~\ref{sec:mechanism} presents the counterfactual, drop-in ablations, and the cross-attention note. Sections~\ref{sec:discussion}~and~\ref{sec:conclusion} close.
